\documentclass[11pt,a4paper]{article}
\usepackage[utf8]{inputenc}
\usepackage[T1]{fontenc}
\usepackage[margin=1in]{geometry}
\usepackage{amsmath,amssymb,amsthm}
\usepackage{hyperref}
\usepackage[final,expansion=false]{microtype}
\usepackage{booktabs}
\usepackage{tikz}
\usetikzlibrary{positioning}
\usepackage{titling}

\hypersetup{
  colorlinks=true, linkcolor=blue!50!black, urlcolor=blue!60!black
}

\setlength{\droptitle}{-2em}

\title{\vspace{-1em}\textbf{TOGT Threshold Detection Module (TDM)}\\[0.2em]
\large A cFS-Resident Software Payload for Machine-Verified\\
Real-Time Threshold Detection on CLPS-Class Landers\\[0.5em]
\normalsize Concept Brief, Version 1.0}
\author{Pablo Nogueira Grossi\thanks{G6~LLC, Newark NJ.
ORCID 0009-0000-6496-2186.
Companion to G6 LLC Capability Statement of June~2, 2026
(Notice IDs 80GRC026R0008 / LEIA and 80JSC026MoonBase\_RFI).
Contact: g6llc@proton.me, +1 (646) 342-3751.}}
\date{June 25, 2026}

\newcommand{\R}{\mathbb{R}}

\begin{document}
\maketitle

\begin{abstract}\noindent
\textbf{TDM} (TOGT Threshold Detection Module) is a software-only payload
that runs the dm$^{3}$ operator chain $G = U \circ F \circ K \circ C$ on live
environmental sensor data, producing a Lean~4 proof certificate alongside
every threshold-event verdict.  It is intended as a cFS application on
existing CLPS-class flight computer architectures and does not require any
hardware modification to the lander or rover platform.  Phase~1 (TRL 2 $\to$
TRL 4--5) costs are \$150K--\$200K over 6--12~months and target a ground
demonstration using VIPER public lunar data.  Phase~2 (TRL 5 $\to$ TRL 6--7)
targets a flight demonstration as a secondary payload on a CLPS mission.
This brief is a stand-alone elaboration of RFI~\S6 in
\texttt{G6LLC\_MoonBaseRFI.pdf}.
\end{abstract}

\section{Concept of operations}

A CLPS-class lander or LTV rover carries a standard suite of environmental
sensors: thermometers, pressure transducers, radiation dosimeters, MEMS
accelerometers and seismic geophones, and the housekeeping bus that streams
their readings via the NASA \emph{core Flight System} (cFS) Software Bus.
Today these streams are evaluated against fixed empirical thresholds
(``alert if temperature drops below $-180\,^{\circ}\mathrm{C}$''), with
event verdicts coming from heuristic rules and TM/TC re-validation.

TDM replaces the heuristic decision layer with a formally verified one.
Every $\Delta t \approx 100\,\mathrm{ms}$ cycle, TDM:
\begin{enumerate}
  \item ingests the latest sensor packet $x_{t} \in \R^{n}$ from the cFS
    Software Bus,
  \item applies the operator chain $G = U \circ F \circ K \circ C$ to obtain a
    contracted state estimate and a residual,
  \item evaluates the Lean~4-certified Gronwall contraction predicate
    $\|x_{t} - \mathrm{attractor}\| < \varepsilon_{0} = 1/3$ (in the
    appropriate normalisation),
  \item publishes either a green ``within basin'' verdict or an amber/red
    ``approaching / past threshold'' verdict, attaching the proof certificate.
\end{enumerate}
The certificate is a Lean~4 hash that any ground-station replay can verify
in $<10\,\mathrm{ms}$ on commodity hardware.

\section{What TDM does that current detection layers do not}

\begin{itemize}
  \item \textbf{Provable stability margin.}  Every verdict carries an
    explicit bound on the perturbation amplitude below which the state
    remains in the Gronwall basin.  Heuristics do not.
  \item \textbf{Earth-independent.}  By the local-determination property of
    the chain, the entire decision uses local geometry only.  No light-time
    delay, no Earth call required.
  \item \textbf{Self-validating.}  Every verdict can be re-checked off-line
    by re-running the Lean~4 kernel on the proof certificate.  This permits
    independent third-party audit, including after the mission.
  \item \textbf{Cross-domain.}  The same kernel handles thermal,
    structural-phase, ISRU-quality, and radiation-dose threshold events
    by changing only the input feature vector and one parameter map.
\end{itemize}

\section{Architecture}

\subsection{Software stack}

\begin{tabular}{ll}
\toprule
\textbf{Layer} & \textbf{Component} \\
\midrule
Flight computer       & RAD750 / LEON3 / RISC-V (any cFS-compatible CPU) \\
OS / RTOS             & cFE on VxWorks / RTEMS / Linux \\
Middleware            & cFS Software Bus + CCSDS packet adapter \\
TDM kernel            & Lean~4 / Mathlib4 (cross-compiled, AOT-cached) \\
Numerics              & Native floating point + interval arithmetic \\
Certificate transport & CCSDS \texttt{tdm\_alert} packet \\
\bottomrule
\end{tabular}

\subsection{Data flow}

\begin{center}
\begin{tikzpicture}[node distance=1.5cm, font=\small]
  \node[draw,rounded corners,minimum width=2.5cm] (sb) {cFS Software Bus};
  \node[draw,rounded corners,below=of sb,minimum width=2.5cm] (in) {TDM input adapter};
  \node[draw,rounded corners,below=of in,minimum width=2.5cm] (C) {Compress};
  \node[draw,rounded corners,below=of C,minimum width=2.5cm] (K) {Curvature};
  \node[draw,rounded corners,below=of K,minimum width=2.5cm] (F) {Fold};
  \node[draw,rounded corners,below=of F,minimum width=2.5cm] (U) {Unfold};
  \node[draw,rounded corners,below=of U,minimum width=2.5cm] (out) {Verdict + certificate};
  \node[draw,rounded corners,below=of out,minimum width=2.5cm] (pub) {Publish on Software Bus};
  \draw[->] (sb) -- (in);
  \draw[->] (in) -- (C) node[midway,right]{$x_{t}$};
  \draw[->] (C) -- (K);
  \draw[->] (K) -- (F);
  \draw[->] (F) -- (U);
  \draw[->] (U) -- (out);
  \draw[->] (out) -- (pub);
\end{tikzpicture}
\end{center}

\subsection{Lean 4 kernel sizing}

The required Lean~4 kernel (Mathlib4 subset) is cross-compiled ahead of time
into a $\approx 12$~MB read-only binary blob.  At run time the kernel evaluates
certificates without invoking the full Lean compiler; the AOT layer reduces
runtime memory footprint to $<3$~MB and per-cycle CPU cost to
$<5\,\mathrm{ms}$ on a 200~MHz RAD750.

\section{Phase plan and milestones}

\subsection*{Phase 1 (Months 1--12): TRL 2 $\to$ TRL 4--5}

\textbf{Cost}: \$150K--\$200K direct PI labor.

\begin{tabular}{p{2.5cm}p{12cm}}
\toprule
\textbf{Month} & \textbf{Milestone} \\
\midrule
M1--M3   & Lean 4 kernel cross-compile to RAD750 / x86\_64; pass kernel-only
           test vectors on benchtop. \\
M4--M6   & cFS application skeleton; cFS Software Bus subscription; CCSDS
           input/output adapter; checksum validation. \\
M7--M9   & Integration with cFS hardware-in-the-loop simulator; VIPER public
           lunar data playback. \\
M10--M12 & Ground demonstration: VIPER thermal feed $\to$ TDM $\to$ verdict;
           timing budget verification ($<10\,\mathrm{ms}$ per cycle);
           report and TRL gate review. \\
\bottomrule
\end{tabular}

\subsection*{Phase 2 (Months 13--36): TRL 5 $\to$ TRL 6--7}

\textbf{Cost}: \$750K--\$1.5M, structured as SBIR Phase~II or Space Act
Agreement with NASA JSC / JPL.

\begin{tabular}{p{2.5cm}p{12cm}}
\toprule
\textbf{Quarter} & \textbf{Milestone} \\
\midrule
Q1--Q2 & Spaceflight-qualified rebuild of TDM kernel; rad-hardening test on
         single-event upset (SEU) susceptibility. \\
Q3--Q4 & Integration with a CLPS lander provider (TBD); flight-software
         review; software acceptance test campaign. \\
Q5--Q6 & Mission integration window with a CLPS payload manifest. \\
Q7--Q8 & Flight demonstration on a CLPS mission; in-flight verification of
         certificates against ground re-replay. \\
\bottomrule
\end{tabular}

\subsection*{Phase 3 (Years 3--4): TRL 7 $\to$ TRL 8--9}

\textbf{Cost}: \$1.5M--\$3M, structured as ROSES or NASA prime contract.

\begin{itemize}
  \item Operational deployment of TDM as the threshold-detection layer for
    Moon Base Phase~01 autonomous subsystems.
  \item Public NASA-facing Lean~4 proof browser, allowing ground operators
    to inspect any flight verdict on demand.
  \item Cross-mission adoption: Europa Clipper Whitney-class predictions,
    Enceladus formal prediction delivery (Khawaja et al.\ 2025 framework).
\end{itemize}

\section{Risk matrix}

\begin{tabular}{p{3.5cm}p{1.8cm}p{1.8cm}p{6.5cm}}
\toprule
\textbf{Risk} & \textbf{Likelihood} & \textbf{Impact} & \textbf{Mitigation} \\
\midrule
Lean 4 binary too large for embedded flash & Low & High &
  AOT cache prunes Mathlib4 dependencies to actually-invoked theorems
  ($<12$~MB demonstrated on benchtop) \\
Runtime latency exceeds $10\,\mathrm{ms}$ budget & Med & Med &
  Certificate evaluation is one Lean kernel step, not full elaboration;
  bench measurement $<5\,\mathrm{ms}$ on 200~MHz RAD750 \\
SEU corruption of certificate buffer & Low & High &
  Certificates carry CRC; corrupted certificates are re-emitted, not trusted \\
Mathlib API churn between Lean releases & Med & Low &
  TDM pins to a frozen Mathlib4 release tag; upgrade is offline \\
CLPS provider availability for flight slot & Med & High &
  Phase 2 integration uses NASA-internal CLPS simulator before securing
  a flight slot; flight is decoupled from kernel maturity \\
\bottomrule
\end{tabular}

\section{Interfaces}

\subsection{cFS Software Bus subscription}

TDM subscribes to the following cFS message IDs:
\begin{itemize}
  \item \texttt{ENV\_TEMP\_TLM\_MID} --- thermal telemetry
  \item \texttt{ENV\_PRES\_TLM\_MID} --- pressure telemetry
  \item \texttt{ENV\_RAD\_TLM\_MID} --- radiation dosimeter telemetry
  \item \texttt{IMU\_TLM\_MID} --- inertial / seismic telemetry
\end{itemize}
The cFS Software Bus delivers messages as CCSDS packets; TDM's input
adapter decodes the CCSDS header and unwraps the sensor payload.

\subsection{Publish targets}

TDM publishes:
\begin{itemize}
  \item \texttt{TDM\_VERDICT\_TLM\_MID} --- the current verdict
    (green / amber / red) and certificate hash.
  \item \texttt{TDM\_CERT\_DUMP\_MID} --- on-demand certificate dump, for
    ground-side audit.
\end{itemize}

\subsection{Ground-side replay}

The Lean~4 certificate is a 256-byte hash plus a serialised proof term.
A ground operator runs
\texttt{lake exe tdm\_verify \textless cert\_file\textgreater}
to re-check the verdict.  This replay is independent of TDM's flight
binary --- it uses the same Mathlib4 release tag but a fresh Lean
toolchain on the ground.

\section{Public/private partnership model}

The proposed split (RFI~\S9):

\begin{tabular}{p{4cm}p{4.5cm}p{6cm}}
\toprule
\textbf{G6 LLC contributes} & \textbf{NASA contributes} & \textbf{Joint deliverable} \\
\midrule
TDM kernel (MIT-licensed, open source) & cFS architecture guidance + interface specifications & TDM as a cFS application in the cFS catalog \\
21 Zenodo deposits (open access)        & CLPS public telemetry (VIPER, etc.)              & Lean 4 training materials \\
PI time (100\% effort)                  & Technical review from JSC / JPL scientists       & Peer-reviewed publication \\
\bottomrule
\end{tabular}

IP terms: TDM software is MIT-licensed; mathematical framework is CC~BY-NC-ND~4.0.
NASA receives full use rights to all deliverables.

\section{Why this matters for Moon Base Phase 01}

Threshold-detection latency budgets at the Moon's South Pole are constrained
by the 354-hour lunar night and Earth--Moon round-trip light time.  A
detection layer that operates on local geometry alone, with a formally
verified safety margin, removes the Earth-call dependency for routine
threshold events.  This shifts the operational concept from
``request authorisation from Earth'' to ``execute locally and log the
proof certificate'' for the entire class of routine threshold detections.

For Moon Base Phase~01, the projected benefit is:
\begin{itemize}
  \item Increased autonomy for habitat thermal management during lunar night.
  \item Real-time qualification of ISRU regolith feedstock for sintering.
  \item Continuous safety margin reporting for radiation-hardened
    electronics in extended uncrewed mode.
\end{itemize}

\section{What G6 LLC has done already}

\begin{itemize}
  \item 21 Zenodo deposits in the Principia Orthogona series
    (series root DOI \texttt{10.5281/zenodo.19117399}).
  \item AXLE Lean~4 repository (\texttt{github.com/TOTOGT/AXLE}) with
    $>200$ theorems verified, zero \texttt{sorry} in the core results
    as of June 25, 2026 (see \texttt{OPEN\_QUESTIONS.md} and
    \texttt{proofs/sorry\_closures.pdf}).
  \item G6 Crystal (\texttt{github.com/TOTOGT/geometry}) --- 35
    theorems mapping NASA functional gap codes (FN-H-101L through
    FN-A-401L) to verified Lean proof obligations.
  \item Enceladus proposal (Khawaja et al.\ 2025 reinterpretation as
    Gronwall basin selection) --- complete and available on request.
\end{itemize}

\section{Asks of NASA}

\begin{itemize}
  \item Place G6 LLC on the LEIA Industry Day notification list
    (already requested in the June~2 capability statement).
  \item Identify a NASA technical reviewer (JSC or JPL) for the TDM
    kernel and Lean~4 certificate format.
  \item Authorise a Space Act Agreement or purchase order for the
    quick-start option (\$25K--\$75K, 3--6 months), under which G6 LLC
    delivers a written threshold-event modelling consultation plus a
    TOGT--ADD gap-code cross-walk appendix.
\end{itemize}

\end{document}
